\section{Rappels Partie 1}

% --- FRAME 1 : RAPPELS RÉGLEMENTATION & PHYSIQUE ---
\begin{frame}{Rappels : Cadre et Principes}
    Avant d'aborder la suite, rappelons les fondamentaux du Niveau 3 :

    \begin{columns}
        \column{0.5\textwidth}
        \begin{block}{Prérogatives N3}
            \begin{itemize}
                \item \textbf{PA60} : Autonome jusqu'à 60m (avec DP).
                \item \textbf{PA40} : Autonome jusqu'à 40m (sans DP).
                \item \textbf{Responsable} de sa planification et de sa sécurité.
            \end{itemize}
        \end{block}

        \column{0.4\textwidth}
        \begin{block}{Physique (Loi de Mariotte)}
            $$ P \times V = \text{Constante} $$
            La pression augmente de \textbf{1 bar tous les 10m}.
            \begin{itemize}
                \item À 20m ($P_{abs}=3b$) : Consommation $\times 3$.
                \item Autonomie divisée par 3 !
            \end{itemize}
        \end{block}
    \end{columns}
    
    \vspace{0.5cm}
    \textbf{Règle d'or en autonomie :} On planifie \textbf{avant} de s'immerger (Air, Temps, Profondeur, Palanquée).
\end{frame}

% --- FRAME 2 : RAPPELS RISQUES ---
\begin{frame}{Rappels : Les Risques Majeurs}
    \begin{alertblock}{Barotraumatismes (Volumes gaz du corps)}
        \textbf{Cause :} Variation de pression (surtout 0-10m).
        \begin{itemize}
            \item \textbf{Surpression Pulmonaire ($\nearrow$)} : Expiration continue obligatoire.
            \item \textbf{Oreilles ($\searrow$)} : Manœuvres douces, ne jamais forcer.
        \end{itemize}
    \end{alertblock}

    \begin{block}{Accident de Désaturation (ADD - Loi de Henry)}
        \textbf{Cause :} Bulles d'azote bloquant la circulation.
        \begin{itemize}
            \item \textbf{Facteurs favorisants} : Profil, Froid, Effort, Age.
            \item \textbf{Prévention} : Respect vitesse remontée, hydratation, pas d'effort après.
            \item \textbf{CAT} : O2 (15L/min) + Évacuation + Eau.
        \end{itemize}
    \end{block}
\end{frame}

% --- FRAME 3 : EXERCICE AVANCÉ ---
\begin{frame}{Exercice de Planification (Niveau N3)}
    \textbf{Scénario :}
    Plongée sur une épave à \textbf{30 mètres}.
    \begin{itemize}
        \item Bloc de \textbf{15 Litres} gonflé à \textbf{210 bars}.
        \item Consommation surface (stress inclus) : \textbf{25 L/min}.
        \item Réserve de sécurité obligatoire : \textbf{50 bars}.
    \end{itemize}

    \vspace{0.3cm}
    \textbf{Calculez votre autonomie au fond selon 3 méthodes :}
    \begin{enumerate}
        \item \textbf{Méthode "Naïve"} : On consomme tout le gaz utilisable au fond (sans penser à la remontée).
        \item \textbf{Méthode "Réaliste"} : On déduit le gaz nécessaire à la remontée (DTR).
        \textit{\small(Pour l'exercice, on estime la consommation de la remontée à \textbf{400 Litres}).}
        \item \textbf{Méthode "Sécurisée"} : On applique la \textbf{Règle des Tiers}.
    \end{enumerate}
\end{frame}

% --- FRAME 4 : SOLUTION ---
% --- SOLUTION PARTIE 1 : SIMPLIFIÉE ---
\begin{frame}{Solution 1 : Méthode "Naïve"}
    \textit{"Je reste au fond tant que je n'ai pas atteint ma réserve."}
    
    \begin{enumerate}
        \item \textbf{Calcul de la Conso Fond ($C_{fond}$)} :
        $$ P_{abs} = 4 \text{ bars} \quad \Rightarrow \quad C_{fond} = 25 \times 4 = \mathbf{100 \text{ L/min}} $$
        
        \item \textbf{Gaz Utilisable :}
        $$ 210b - 50b (\text{Réserve}) = 160b $$
        $$ 160b \times 15L = \mathbf{2400 \text{ Litres}} $$
        
        \item \textbf{Autonomie :}
        $$ \text{Temps} = \frac{2400}{100} = \mathbf{24 \text{ minutes}} $$
    \end{enumerate}
    
    \begin{alertblock}{Danger}
        Si vous restez 24 min, vous entamez la remontée avec 50 bars.
        Vous arriverez en surface \textbf{dans la réserve} (car vous allez consommer en remontant). \textbf{C'est risqué.}
    \end{alertblock}
\end{frame}

% --- SOLUTION PARTIE 2 : AVEC DTR ---
\begin{frame}{Solution 2 : Intégration de la DTR}
    \textit{"Je garde du gaz pour remonter sereinement."}
    
    \begin{itemize}
        \item \textbf{Gaz Utilisable Total} : 2400 Litres (cf. calcul précédent).
        \item \textbf{Besoin Remontée (DTR)} : 400 Litres (Donnée exercice).
    \end{itemize}
    
    $$ \text{Gaz Dispo Fond} = 2400 - 400 = \mathbf{2000 \text{ Litres}} $$
    
    \textbf{Nouvelle Autonomie :}
    $$ \frac{2000}{100} = \mathbf{20 \text{ minutes}} $$
    
    \begin{block}{Constat}
        On perd 4 minutes de temps fond, mais on sort de l'eau avec la réserve intacte (50b). C'est la procédure standard minimale.
    \end{block}
\end{frame}

% --- SOLUTION PARTIE 3 : TIERS ---
\begin{frame}{Solution 3 : La Règle des Tiers}
    \textit{"Je consomme 1/3 à l'aller, 1/3 au retour, 1/3 sécurité."}
    \begin{columns}
        \column{0.4\textwidth}
        \begin{enumerate}
            \item \textbf{Calcul du Tiers (en pression) :}
            $$ \frac{210 \text{ bars}}{3} = \mathbf{70 \text{ bars}} $$
            
            \item \textbf{Volume disponible (1/3) :}
            $$ 70 \text{ bars} \times 15L = \mathbf{1050 \text{ Litres}} $$
            
            \item \textbf{Autonomie (Demi-tour) :}
            $$ \frac{1050}{100} \approx \mathbf{10 \text{ minutes}} $$
        \end{enumerate}

        \column{0.5\textwidth}
        \begin{exampleblock}{Comparaison Finale à 30m}
            \begin{itemize}
                \item Calcul simpliste : \textbf{24 min} (Dangereux)
                \item Calcul DTR : \textbf{20 min} (Standard)
                \item Règle des Tiers : \textbf{10 min} (Très sécurisé)
            \end{itemize}
        \end{exampleblock}
    \end{columns}
\end{frame}